\section{Model}
\label{sec:model}

The model is an individual-based, discrete-time stochastic process on a lattice, with one step per week. Each
individual has a position, an age, a sex, a sparse affinity memory and a \emph{social state} $s_i\in[0,1]$, its
capacity for normal social and reproductive behavior. All rules are local: no individual knows the population size. The \emph{core model}
(lattice, affinity, movement with attraction, crowding damage, reproduction and senescent mortality) is extended by
four rules, R1, R2, AV and R7, each with a neutral value that switches it off and recovers the core model
(Fig.~\ref{fig:model-schematic}, Table~\ref{tab:params}; ODD in the Supplemental Material).

\begin{figure}[t]
  \centering
  \includegraphics[width=\columnwidth]{model_schematic_pre.pdf}
  \caption{Schematic of the model.
  (a)~A focal agent (ringed) chooses a cell of its Moore neighborhood (dashed) with the weight of
  Eq.~\eqref{eq:model-move}, whose three factors are shown below the lattice; arrow width is proportional to the
  weight. Deteriorated individuals gather in pools ($m>m_c$) or occupy refuges (hatched, capacity $c_R=2$).
  (b)~Illustrative trajectories of $s$ versus age, from Eqs.~\eqref{eq:model-social} and~\eqref{eq:model-learning}
  with fixed environments (not a simulation). A pup starts at $w\,s_{\rm mother}$ and gains competence only inside
  the window $a\le a_w$ (R1), at a rate set by its neighbors (R2); afterwards crowding can only lower $s$.}
  \label{fig:model-schematic}
\end{figure}

\subsection{Core model}
\label{sec:model-baseline}

\paragraph*{Space, affinity and movement.}
Individuals live on an $L\times L$ lattice $G$ ($L=30$, $|G|=900$) with hard walls and no occupancy limit; $m_i$ is
the occupancy of the cell of $i$, itself included. Pairs at Chebyshev distance $\le1$ reinforce their affinity,
$F_{ij}\leftarrow(1-\eta)F_{ij}+\eta$, and the others decay, $F_{ij}\leftarrow(1-\delta)F_{ij}$: the memory,
$1/\delta=100$ weeks, is long and built only by proximity. Every week all individuals move simultaneously to one of
the nine cells $y$ of their Moore neighborhood, own cell included, with probability proportional to
\begin{equation}
  \begin{split}
  W_i(y)={}&\bigl[1+\alpha A_i(y)\bigr]\,e^{-\kappa(1-s_i)^{p}n_y}\\
           &\times\bigl[1+\pi(1-b_i)^{q}\,\mathbf 1_{\mathcal R}(y)\bigr],
  \end{split}
  \label{eq:model-move}
\end{equation}
where $A_i(y)=\sum_{j\in y}F_{ij}$, $n_y$ is the number of \emph{other} individuals in $y$, $\mathcal R$ is the set
of non-full refuges (refuge cells with $n_y<c_R$), and $\mathbf 1_{\mathcal R}(y)=1$ if $y$ is a non-full refuge and 0
otherwise; full refuges cannot be entered (R7, below). The core model has only the first factor, the attraction $\alpha$ towards acquaintances; the
other two are the rules AV and R7.

\paragraph*{Social state.}
After moving,
\begin{equation}
  s_i\leftarrow\min\bigl\{1,\max\bigl[0,\;s_i+\varrho_i-\gamma\,(m_i-m_c)_+\bigr]\bigr\},
  \label{eq:model-social}
\end{equation}
where $\gamma(m_i-m_c)_+$ is the damage inflicted by each occupant in excess of the community size $m_c$, and in the
core model recovery is a constant, $\varrho_i=r$. Deterioration is local, density-triggered and, being convex in the
occupancy, acts as a switch.

\paragraph*{Demography.}
A reproductive female ($a_{\min}\le a\le a_{\rm rep}$) sharing her cell with a reproductive male conceives with
probability $p_{\rm conc}=p_0\,s_f\,s_m$; each pup of the litter (mean 6.55) survives birth with probability $q_0\,s_f$
and is born with $s_{\rm nb}=w\,s_f+(1-w)\,s_0$ (full inheritance, $w=1$, in the core model). Death is senescent, with
weekly probability $\{1+\exp[-k(a-a_{\max})]\}^{-1}$, independent of $s$ (an optional social advance of
senescence, $\mu_s$, is zero in every model): the beautiful ones were ``physically
healthy''~\cite{Calhoun1973}, so the absence of social deaths (A4) and the senescent decline (O8) hold by
construction. The long-lived demography of the candidates ($a_{\max}=120$, $a_{\rm rep}=80$ weeks) matches the menopause at about 560
days and the longevity of about 800 days in Universe~25~\cite{Calhoun1973}. With no mortality before senescence the
population sits on a plateau once natality stops. Initially $200+200$ individuals of age 4 weeks with $s=1$ are placed at random.

\paragraph*{Controls.}
A rule is switched off by its neutral value: $a_w=\infty$ (R1); $\lambda=0$ with $w=1$ and individual recovery at
$r=0.01$ (R2); $\kappa=0$ (AV); no refuges (R7). The \emph{short-lived control} is the core model with all four rules
switched off, attraction $\alpha=1.4$ and a short life span, $a_{\max}=60$ and $a_{\rm rep}=50$ weeks; the
\emph{long-lived control} is the same with $a_{\max}=120$ and $a_{\rm rep}=80$ (Table~\ref{tab:params}). The null
models are defined relative to \textsf{C1} in Sec.~\ref{sec:methods}.

\subsection{Minimal additional rules}
\label{sec:model-rules}

With the long-lived demography and $\alpha=4$, crowding alone already produces the envelope of the collapse, but not
coexisting classes of deteriorated animals (Sec.~\ref{sec:results}). Each rule is motivated by an observation of the
original report~\cite{Calhoun1973}; R7 (five parameters and a hard capacity) was designed after a candidate without
refuges had failed to produce a persistent isolated class.

\paragraph*{R1: juvenile plasticity window ($a_w$).}
Recovery is possible only for $a\le a_w=12$ weeks ($\varrho_i=0$ otherwise), while crowding damage acts at every age: an adult
is a ratchet, $\dot s_i\le0$, and recovery becomes a lineage process. Since $p_{\rm conc}\propto s_fs_m$, an adult with
$s\simeq0$ can neither recover nor reproduce, whatever its partner. \emph{Motivation:} mice moved in mid-phase~D to
uncrowded universes, even with partners matured without crowding, never recovered reproductive behavior (Marsden's
work~\cite{Calhoun1973,Marsden1972}); the failure of phase-D transplants (O12) is thus a
consistency check of R1. \emph{Mechanism:} refuges are pockets of low density in which, without R1, deteriorated
animals recover and breed again (natality resumes and the primary outcome falls to 0.21), whereas without refuges the
collapse stays irreversible in most runs even without R1 (O6 and O7 in 0.85 of 40 runs). R1 is what lets a persistent isolated class
coexist with an irreversible collapse.

\paragraph*{R2: social learning ($\lambda$, with partial inheritance $w$).}
Recovery is \emph{learned} from the neighbors,
\begin{equation}
  \varrho_i=r\bigl[(1-\lambda)+\lambda\,\bar s_{V(i)}\bigr],
  \label{eq:model-learning}
\end{equation}
with $\bar s_{V(i)}$ the mean state of the other individuals of the $3\times3$ block centred on $i$ (0 if none).
With R1 it acts only on juveniles, which start with a fifth of their mother's competence ($w=0.2$, $s_0=0$) and can
gain at most $r\,a_w=1.2$ from competent neighbors. For $\lambda=1$, $s=0$ is absorbing for the population, not only
for the individual: without competent adults pups have nobody to learn from.
\emph{Motivation:} ``by midway in phase~C essentially all young were prematurely rejected by their
mothers''~\cite{Calhoun1973}. Pups are born at $s_{\rm nb}\le0.2$, the patterns' threshold for ``deteriorated,'' so
part of the failure of crowded-born cohorts is definitional (Sec.~\ref{sec:results}); the essential part of R2 is
learning, not inheritance ($w=0$ with $r=0.15$ passes in all 120 runs).

\paragraph*{AV: crowd aversion of the deteriorated ($\kappa$).}
The factor $e^{-\kappa(1-s_i)^{p}n_y}$ of Eq.~\eqref{eq:model-move}, with $p=4$, acts only for $s\lesssim0.3$.
Facing a cell of $n$ acquaintances ($A\simeq n$), a deteriorated individual maximizes $(1+\alpha n)e^{-\kappa n}$ at
$n^*=1/\kappa-1/\alpha\approx6.4\approx0.8\,m_c$ and avoids cells of strangers, so the deteriorated split between
pools, where damage continues, and nearly empty cells. \emph{Motivation:} withdrawn males pooled and wounded each other; the beautiful ones ``never engaged in
fighting''~\cite{Calhoun1973}. Without refuges AV alone generates the isolated class; with refuges it is not needed for
the primary outcome and instead shapes the use of space, \emph{reducing} the empty fraction of cells at the peak
(Sec.~\ref{sec:results}).

\paragraph*{R7: refuges of hard capacity ($\pi$).}
A fixed random fraction $f_R$ of the cells are refuges that hold at most $c_R=2$ individuals (a hard capacity:
residents keep their place, rejected entrants return). The deteriorated prefer refuges through the factor $1+\pi(1-b_i)^{q}$ of Eq.~\eqref{eq:model-move}, with $b_i=s_i$. The total capacity,
$f_Rc_R|G|=540$ places or about 20\% of the adults at the peak, bounds the persistently isolated fraction. The hard
capacity remedies the stampedes of synchronous moves (Sec.~\ref{sec:results}). \emph{Motivation:} withdrawn females retreated to the high, less preferred nest
boxes~\cite{Calhoun1973}, the analogue of the refuges if a cell is a nest box; but $c_R=2$ is the isolation threshold
$m\le2$ of O13, a design choice that makes the persistence of the isolated class (O13p) hold by construction.

\paragraph*{Candidates and time scales.}
\label{sec:model-candidates}
Rules and parameters were chosen to maximize the number of patterns reproduced with the fewest rules, not by fitting a curve,
and with the evaluator later used to report them, so calibration is not validation (Sec.~\ref{sec:methods}). The reference candidate \textsf{C1} has the four rules. \textsf{C1$'$} omits AV
and is the minimal set for the preregistered primary outcome O17p (each single ablation gives O17p $=0$ in 40 runs),
but only at the evaluation thresholds: it passes in 0.835 [0.78, 0.88] of the runs, against 0.995 for \textsf{C1}.
Deterioration takes about one
generation: the animals past the window lose half of their competence $t_{\rm det}\approx10$ weeks after crowding sets
in, and $\Gamma=t_{\rm det}/T_{\rm gen}\approx1.1$ (0.7 with the mean age of the mothers at birth as $T_{\rm gen}$), a
heuristic of scales against $\Gamma\approx3.5$ for Universe~25 (our estimate, with a different operational
definition) that does not set the depth of the boom.

\begin{table}[t]
  \caption{Parameters of the reference candidate \textsf{C1} and of the short-lived control (SL); \textsf{C1$'$}
  is \textsf{C1} with $\kappa=0$, and the long-lived control is SL with $a_{\max}=120$ and $a_{\rm rep}=80$. Time in weeks; a dash means that the rule is inactive. The $N_0$ initial individuals
  have age 4 weeks and $s=1$. Presets \texttt{calhoun\_C1} and \texttt{calhoun\_C1p} of the \texttt{u25} package.}
  \label{tab:params}
  \begin{ruledtabular}
  \begin{tabular}{lcclcc}
    Symbol & \textsf{C1} & SL & Symbol & \textsf{C1} & SL \\
    \hline
    $L$ & 30 & 30 & $a_w$ (R1) & 12 & $\infty$ \\
    $N_0$ & 400 & 400 & $\lambda$ (R2) & 1 & 0 \\
    $\alpha$ & 4 & 1.4 & $w,\ s_0$ (R2) & 0.2, 0 & 1, -- \\
    $\eta,\ \delta$ & 0.2, 0.01 & 0.2, 0.01 & $\kappa,\ p$ (AV) & 0.15, 4 & -- \\
    $m_c,\ \gamma$ & 8, 0.1 & 8, 0.1 & $f_R,\ c_R$ (R7) & 0.3, 2 & -- \\
    $r$ & 0.1 & 0.01 & $\pi,\ q,\ b_i$ (R7) & 100, 4, $s_i$ & -- \\
    $p_0,\ q_0$ & 0.3, 0.95 & 0.3, 0.95 & $a_{\min},\ a_{\rm rep}$ & 6, 80 & 6, 50 \\
    $\mu_s$ & 0 & 0 & $a_{\max},\ k$ & 120, 0.15 & 60, 0.15 \\
  \end{tabular}
  \end{ruledtabular}
\end{table}
